% Work/MAT — přepracovaná příprava hodiny 14: Vennovy diagramy.
% Návaznost: 12 Co je množina; 13 Podmnožiny a čtyři množinové operace.
% V další hodině: vlastnosti množinových operací a jejich zdůvodnění.
% Učebnice autora: ucebnice/01_Zakladni_poznatky/Mnoziny/Množiny.tex,
%   oddíly Ilustrace, Množinové operace, Vennovy diagramy a slovní úlohy.
% Realisticky.cz: 010304 Vennovy diagramy — práce se čtyřmi a osmi poli,
%   zásada umístit prvek právě jednou, rozlišení obsazeného a prázdného pole,
%   rozbor kombinací příslušnosti. Metodická inspirace; níže vlastní úlohy.
% https://www.realisticky.cz/ucebnice/01%20Matematika%20S%C5%A0/01%20Z%C3%A1kladn%C3%AD%20poznatky/03%20Mno%C5%BEiny/04%20Vennovy%20diagramy.pdf
% 010304 Příklady: čtení diagramu a tvorba třímnožinového diagramu.
% https://www.realisticky.cz/ucebnice/01%20Matematika%20S%C5%A0/01%20Z%C3%A1kladn%C3%AD%20poznatky/03%20Mno%C5%BEiny/04%20Vennovy%20diagramy/p%C5%99%C3%ADklady.pdf
% Všechna následující zadání a diagramy jsou původní didaktické zpracování.
% Šablona: ../../sablony/pripravy/priprava.sty (neměnit).
\documentclass[11pt,a4paper]{article}
% Spolecny styl pro ucitelske pripravy MAT (Work/MAT)
% Verze 1.1, 2026-10-23
\NeedsTeXFormat{LaTeX2e}
\ProvidesPackage{priprava}[2026/10/23 v1.1 Sablona ucitelskych priprav MAT]

% Zakladni typografie a cestina
\RequirePackage[T1]{fontenc}
\RequirePackage[utf8]{inputenc}
\RequirePackage[czech]{babel}
\RequirePackage{lmodern}
\RequirePackage{microtype}
\RequirePackage{amsmath,amssymb}
\RequirePackage{xcolor}
\RequirePackage{enumitem}
\RequirePackage[most]{tcolorbox}
\RequirePackage{tabularx,booktabs,array}
\RequirePackage{fancyhdr}
\RequirePackage{needspace}
\RequirePackage{geometry}
\RequirePackage[hidelinks]{hyperref}

\geometry{a4paper,top=17mm,bottom=18mm,left=23mm,right=23mm}
\definecolor{notegray}{gray}{0.39}
\setlength{\parindent}{0pt}
\setlength{\parskip}{2.5pt plus 1pt minus 0.5pt}
\setlist[itemize]{topsep=2pt,itemsep=1pt,parsep=0pt,leftmargin=1.45em}
\setlist[enumerate]{topsep=2pt,itemsep=2pt,parsep=0pt,leftmargin=1.65em}
\setcounter{secnumdepth}{0}
\renewcommand{\arraystretch}{1.12}

% Zahlavi: vlevo rocnik | globalni tema, vpravo cislo: nazev hodiny
% Zapati zustava zcela prazdne.
\pagestyle{fancy}
\fancyhf{}
\renewcommand{\headrulewidth}{0.35pt}
\renewcommand{\footrulewidth}{0pt}
\setlength{\headheight}{14pt}
\newcommand{\nastavPripravu}[4]{%
  \fancyhead[L]{\small #1\enspace\textbar\enspace #2}%
  \fancyhead[R]{\small #3: #4}%
}

% Minimalisticke cernobile boxy. Kazda definice zvlast.
\tcbset{commonlesson/.style={enhanced,breakable=false,colback=white,colframe=black,
  colbacktitle=white,coltitle=black,boxrule=0.45pt,
  arc=0.8mm,outer arc=0.8mm,left=3mm,right=3mm,top=1.6mm,bottom=1.6mm,
  toptitle=1mm,bottomtitle=0.7mm,fonttitle=\bfseries,
  attach boxed title to top left={xshift=2.5mm,yshift=-1.15mm},
  boxed title style={colback=white,colframe=white,boxrule=0pt,left=1mm,right=1mm}}}
\newtcolorbox{zadaniBox}[1][]{commonlesson,title={Zadání},#1}
\newtcolorbox{prikladBox}[1][]{commonlesson,title={Příklad},#1}
\newtcolorbox{definiceBox}[1][]{commonlesson,title={Definice},#1}

% Spolecny box pro zadani a jeho reseni, bez mezery mezi castmi.
% Pouziti:
%   \begin{prikladReseniBox}
%     Zadani ulohy.
%     \tcblower
%     \begin{align*} ... \end{align*}
%   \end{prikladReseniBox}
% Spodni cast ma jemny sedy podklad a delici linku; lze v ni sazet
% vice radku matematiky. Obe casti maji jen jeden spolecny ramecek.
\tcbset{solutionlesson/.style={
  commonlesson,
  enhanced jigsaw,
  breakable,
  colbacklower=black!3,
  segmentation style={solid,draw=black!25,line width=0.35pt},
  fontlower=\small\color{notegray},
  before lower={\noindent\textbf{Řešení:}\enspace},
  leftlower=3mm,rightlower=3mm
}}
\newtcolorbox{prikladReseniBox}[1][]{solutionlesson,title={Příklad},#1}
\newtcolorbox{zadaniReseniBox}[1][]{solutionlesson,title={Zadání},#1}

% Pomucky pro autora pripravy: sedive poznamky; \reseni pro starsi pripravy
\newcommand{\poznamka}[1]{\par\begingroup\small\color{notegray}\itshape
  \textnormal{\textbf{Poznámka pro učitele.}}\enspace #1\par\endgroup}
\newcommand{\reseni}[1]{\par\begingroup\small\color{notegray}
  \textbf{Očekávané řešení:}\enspace #1\par\endgroup}
\newcommand{\blok}[1]{\needspace{5\baselineskip}\subsection*{#1}}
\newcommand{\vyrok}[1]{\uv{#1}}

\usepackage{tikz}
\usetikzlibrary{calc,patterns}
\nastavPripravu{1. ročník}{Množiny}{14}{Vennovy diagramy}

% Diagram dvou množin, měřítko vnějšího volání. Argument = obsah výplně.
\newcommand{\vennTwo}[1]{%
\begin{tikzpicture}[scale=.58,baseline=(current bounding box.center)]
  \path[use as bounding box] (-3.25,-2.12) rectangle (3.25,2.12);
  #1
  \draw[thick] (-3.05,-1.86) rectangle (3.05,1.88);
  \draw[thick] (-.9,0) circle (1.45);
  \draw[thick] (.9,0) circle (1.45);
  \node at (-1.72,1.53) {$A$};
  \node at (1.72,1.53) {$B$};
  \node[anchor=north east] at (2.97,1.82) {$U$};
\end{tikzpicture}}
\newcommand{\shadeIntersection}{\begin{scope}\clip (-.9,0) circle (1.45);\fill[black!25] (.9,0) circle (1.45);\end{scope}}
\newcommand{\shadeUnion}{\fill[black!25] (-.9,0) circle (1.45);\fill[black!25] (.9,0) circle (1.45);}
\newcommand{\shadeDifference}{\begin{scope}\clip (-.9,0) circle (1.45);\fill[black!25] (-3,-1.8) rectangle (3,1.8);\fill[white] (.9,0) circle (1.45);\end{scope}}
\newcommand{\shadeComplement}{\fill[black!25] (-3.05,-1.86) rectangle (3.05,1.88);\fill[white] (-.9,0) circle (1.45);}
\begin{document}

\blok{Proč kreslit množiny?}
Po minulé hodině umíme spočítat $A\cap B$, $A\cup B$ i $A\setminus B$.
Vennův diagram nám umožní \textbf{vidět, které podmínky platí zároveň}.
Pravidlo je jednoduché: každý uvažovaný prvek umístíme do \textbf{právě
jedné} oblasti. Prvky mimo obě množiny nesmíme zapomenout.

\begin{definiceBox}
\textbf{Základní množina} $U$ je množina všech prvků uvažovaných v úloze.
Ve Vennově diagramu ji obvykle znázorníme obdélníkem; množiny
$A,B,C\subseteq U$ zakreslíme uzavřenými křivkami, zpravidla kružnicemi.
\end{definiceBox}
\begin{definiceBox}
\textbf{Vennův diagram} rozděluje základní množinu do oblastí
podle příslušnosti jednotlivých prvků ke zvoleným množinám.
Pro dvě množiny rozlišujeme \textbf{čtyři} možné oblasti, pro tři
\textbf{osm}, včetně oblasti mimo všechny zakreslené množiny.
\end{definiceBox}

\poznamka{Nejdříve se zeptej, do které oblasti patří věc s oběma
vlastnostmi; teprve potom kresli. Geometrická plocha oblasti nevyjadřuje
počet prvků. Oblast může být prázdná, i když je v obrázku nakreslena.}

\blok{1. Dvě množiny: čtyři různé odpovědi}
\begin{prikladReseniBox}[title={P1 — jádro: roztřiď všechna čísla}]
Nechť $U=\{1,2,\ldots,12\}$. Množina $A$ obsahuje \emph{sudá čísla}
a množina $B$ \emph{násobky tří}. Rozmísti prvky $U$ do diagramu.
Potom napiš, kolik prvků je v každé ze čtyř oblastí.
\tcblower
$A=\{2,4,6,8,10,12\}$ a $B=\{3,6,9,12\}$.
Nejprve zapíšeme společné prvky $6,12$, potom zbývající prvky každé
množiny, nakonec čísla mimo obě množiny.
\begin{center}
\begin{tikzpicture}[scale=.88]
  \draw[thick] (-3.45,-2.08) rectangle (3.45,2.05);
  \draw[thick] (-1,0) circle (1.55);
  \draw[thick] (1,0) circle (1.55);
  \node at (-1.9,1.6) {$A$}; \node at (1.9,1.6) {$B$};
  \node[anchor=north east] at (3.35,1.95) {$U$};
  \node[font=\small] at (-1.70,0) {$2,4,8,10$};
  \node at (0,0) {$6,12$};
  \node at (1.77,0) {$3,9$};
  \node at (0,-1.82) {$1,5,7,11$};
\end{tikzpicture}
\end{center}
\textbf{Jen v $A$:} $\{2,4,8,10\}$ (4 prvky); \textbf{v obou:}
$\{6,12\}$ (2); \textbf{jen v $B$:} $\{3,9\}$ (2);
\textbf{mimo obě:} $\{1,5,7,11\}$ (4).
Kontrola úplnosti: $4+2+2+4=12=|U|$.
\end{prikladReseniBox}

\poznamka{Zeptej se: Kam umístíme $12$? Proč jej nenapíšeme do dvou polí?
Je číslo $7$ uvnitř obdélníku, přestože není v žádném kruhu?
Kontrolu součtem počtů prvků použij až po samostatném pokusu.}

\blok{2. Co vlastně vybarvujeme?}
Obrázek funguje i bez konkrétních čísel: \textbf{šrafovaná či šedá
oblast} odpovídá všem prvkům splňujícím danou podmínku.
Podle zápisu z hodiny 13 vyznačujeme například tyto čtyři oblasti:

\begin{center}
\begin{tabularx}{\linewidth}{@{}XX@{}}
\centering\vennTwo{\shadeIntersection} & \centering\vennTwo{\shadeUnion}\tabularnewline[-.2em]
\centering $A\cap B$ — obě podmínky & \centering $A\cup B$ — alespoň jedna\tabularnewline[.55em]
\centering\vennTwo{\shadeDifference} & \centering\vennTwo{\shadeComplement}\tabularnewline[-.2em]
\centering $A\setminus B$ — v $A$, ne v $B$ & \centering $A'_U=U\setminus A$ — mimo $A$\tabularnewline
\end{tabularx}
\end{center}

\poznamka{Ve výkladu nech žáky nejprve ukázat požadovanou oblast
prstem. Obzvlášť u doplňku připomeň, že vybarvujeme i část druhého
kruhu, pokud leží mimo $A$. U rozdílu rozhoduje pořadí $A\setminus B$.}

\begin{prikladReseniBox}[title={P2 — jádro: čtení výsledků z P1}]
Z diagramu P1 bez nového počítání urči
$A\cap B$, $A\cup B$, $A\setminus B$ a $A'_U$.
\tcblower
$A\cap B=\{6,12\}$;
$A\cup B=\{2,3,4,6,8,9,10,12\}$;
$A\setminus B=\{2,4,8,10\}$;
$A'_U=\{1,3,5,7,9,11\}$.
Doplněk bereme uvnitř $U$: čísla $3,9$ tedy do $A'_U$ patří.
\end{prikladReseniBox}

\blok{3. Tři množiny: osm oblastí}
Nová množina $C$ musí v obrázku procházet všemi dosavadními oblastmi,
aby vzniklo všech osm kombinací příslušnosti. Přidáním třetí kružnice
na správné místo se původní čtyři oblasti rozdělí na dvě části.

\begin{prikladReseniBox}[title={P3 — jádro: společné a výlučné části}]
V $U=\{1,2,\ldots,16\}$ nechť $A$ jsou sudá čísla, $B$ násobky tří
a $C$ čísla menší než $9$. Umísti všechny prvky do diagramu.
Které číslo patří do všech tří množin? Která patří právě do dvou?
\tcblower
$A=\{2,4,6,8,10,12,14,16\}$,
$B=\{3,6,9,12,15\}$,
$C=\{1,2,3,4,5,6,7,8\}$.
Nejprve umístíme $6$ do středu; poté čísla splňující přesně dvě
podmínky a až nakonec čísla splňující jednu nebo žádnou.
\begin{center}
\begin{tikzpicture}[scale=.83]
  \draw[thick] (-3.15,-2.65) rectangle (3.15,2.2);
  \draw[thick] (-.95,.45) circle (1.55);
  \draw[thick] (.95,.45) circle (1.55);
  \draw[thick] (0,-.7) circle (1.55);
  \node at (-1.9,1.65) {$A$}; \node at (1.9,1.65) {$B$};
  \node at (0,-2.37) {$C$}; \node[anchor=north east] at (3.08,2.13) {$U$};
  \node[font=\scriptsize] at (-1.65,.62) {$10,14,16$};
  \node[font=\small] at (1.72,.65) {$9,15$};
  \node at (0,1.26) {$12$};
  \node[font=\small] at (-.86,-.75) {$2,4,8$};
  \node at (.91,-.72) {$3$};
  \node at (0,-.14) {$6$};
  \node[font=\small] at (0,-1.72) {$1,5,7$};
  \node[font=\small] at (2.62,-1.95) {$11,13$};
\end{tikzpicture}
\end{center}
Ve všech třech množinách leží $\{6\}$. Právě do dvou patří
$\{2,3,4,8,12\}$. Ověření: osm oblastí obsahuje
$3+2+1+3+1+1+3+2=16$ čísel.
\end{prikladReseniBox}

\poznamka{Nepřeskakuj společný střed! Při osmi oblastech se často
zapomíná na oblast mimo všechny tři množiny. Není třeba zavádět nový
symbol pro \uv{právě dvě}; postačí slovní popis.}

\blok{4. Samostatné procvičení — jádro}
\begin{zadaniReseniBox}[title={S1 — z diagramu k matematickému zápisu}]
\begin{enumerate}[label=\alph*)]
\item V P1 urči množinu čísel, která jsou v $B$, ale ne v $A$.
Která čísla nepatří ani do $A$, ani do $B$?
\item Z P3 vypiš prvky $A\cap C$, $B\setminus C$ a
$(A\cap B)\setminus C$. U každého výsledku vysvětli oblast slovně.
\item Pro $U=\{a,b,c,d,e,f\}$, $A=\{a,c,e\}$, $B=\{b,c\}$
a $C=\{c,d,e\}$ nakresli diagram tří množin a rozmísti všech šest prvků.
Které dvě oblasti uvnitř kruhů zůstanou prázdné?
\end{enumerate}
\tcblower
\begin{enumerate}[label=\alph*)]
\item $B\setminus A=\{3,9\}$; mimo oba kruhy jsou $\{1,5,7,11\}$.
\item $A\cap C=\{2,4,6,8\}$; $B\setminus C=\{9,12,15\}$;
$(A\cap B)\setminus C=\{12\}$. Poslední výsledek
leží v překryvu $A$ a $B$, ale mimo $C$.
\item Jen v $A$: $a$; jen v $B$: $b$; jen v $C$: $d$;
$A\cap C$ bez $B$: $e$; ve všech třech: $c$; mimo všechny: $f$.
Prázdné jsou $A\cap B$ bez $C$ a $B\cap C$ bez $A$.
\end{enumerate}
\end{zadaniReseniBox}

\poznamka{Při pomalejším tempu dokonči S1 c) doma. U S1 b) dbej,
aby žák nejdříve označil oblast v diagramu a až potom četl čísla.
Nenech čas utéct u samotného kreslení tří kružnic: důležitá je jejich
správná poloha a přesné umístění prvků, nikoli estetická dokonalost.}

\blok{5. Volitelné rozšíření: vybarvi složené výrazy}
\begin{zadaniReseniBox}[title={R1 — dva různé postupy, tentýž diagram}]
V prázdném diagramu tří množin vyznač
\begin{enumerate}[label=\alph*)]
\item $(A\cup B)\setminus C$;
\item $(A\cap B)\cup C$;
\item $A\setminus(B\cup C)$.
\end{enumerate}
Poté u každého výrazu vypiš odpovídající prvky z P3.
\tcblower
\begin{enumerate}[label=\alph*)]
\item Vybarvi oblasti jen $A$, jen $B$ a $(A\cap B)$ mimo $C$:
$\{9,10,12,14,15,16\}$.
\item Vybarvi celý kruh $C$ a navíc $(A\cap B)$ mimo $C$:
$\{1,2,3,4,5,6,7,8,12\}$.
\item Jen $A$ a současně mimo $B$ i $C$:
$\{10,14,16\}$.
\end{enumerate}
\end{zadaniReseniBox}

\blok{6. Domácí procvičení}
\begin{zadaniReseniBox}[title={D1 — nakresli oba diagramy a ověř úplnost}]
\begin{enumerate}[label=\alph*)]
\item Pro $U=\{0,1,\ldots,8\}$, $A=\{0,2,4,6,8\}$
a $B=\{1,2,3,6\}$ rozděl všechna čísla do čtyř oblastí.
Pak urči $A\cap B$ a $A\setminus B$.
\item Pro $U=\{1,\ldots,9\}$,
$A=\{1,2,3,4,5\}$, $B=\{3,4,6,7\}$,
$C=\{2,4,6,8\}$ rozmísti čísla do osmi oblastí.
Urči prvky patřící právě do dvou množin.
\item V úloze b) ověř součtem počtů prvků, že žádné číslo nechybí.
\end{enumerate}
\tcblower
\begin{enumerate}[label=\alph*)]
\item Jen $A$: $\{0,4,8\}$; obě: $\{2,6\}$;
jen $B$: $\{1,3\}$; mimo obě: $\{5,7\}$.
$A\cap B=\{2,6\}$, $A\setminus B=\{0,4,8\}$.
\item Jen $A$: $\{1,5\}$; jen $B$: $\{7\}$;
jen $C$: $\{8\}$; $A,B$ bez $C$: $\{3\}$;
$A,C$ bez $B$: $\{2\}$; $B,C$ bez $A$: $\{6\}$;
$A,B,C$: $\{4\}$; mimo všechny: $\{9\}$.
Právě dvě: $\{2,3,6\}$.
\item Součet počtů je $2+1+1+1+1+1+1+1=9=|U|$.
\end{enumerate}
\end{zadaniReseniBox}


\clearpage
\blok{Příloha: prázdné diagramy k promítání nebo tisku}
\poznamka{\textbf{Jádro 45 minut včetně organizace:}
3 min motivace; 10 min P1; 8 min čtyři vyšrafované oblasti;
7 min P2; 10 min P3; 5 min S1 a), b); 2 min shrnutí a zadání domácí práce.
S1 c) a R1 jsou volitelné. Toto je orientační rozvrh, nikoli povinnost
probrat celý připravený dokument.}

\textbf{Diagram I — dvě množiny (P1, P2, D1 a).}
Doplň do diagramu prvky, zvýrazni požadované oblasti a ověř součet jejich počtů.
\begin{center}
\begin{tikzpicture}[scale=1.25]
  \draw[thick] (-3.35,-1.95) rectangle (3.35,2.05);
  \draw[thick] (-.95,0) circle (1.48);
  \draw[thick] (.95,0) circle (1.48);
  \node at (-2,1.6) {$A$}; \node at (2,1.6) {$B$};
  \node[anchor=north east] at (3.2,1.9) {$U$};
\end{tikzpicture}
\end{center}

\textbf{Diagram II — tři množiny (P3, S1 c), R1, D1 b).}
Začni společnou částí všech tří množin. Označ všech osm oblastí včetně
části mimo všechny kruhy, i když je některá prázdná.
\begin{center}
\begin{tikzpicture}[scale=1.12]
  \draw[thick] (-3.15,-2.65) rectangle (3.15,2.2);
  \draw[thick] (-.95,.45) circle (1.55);
  \draw[thick] (.95,.45) circle (1.55);
  \draw[thick] (0,-.7) circle (1.55);
  \node at (-2.23,1.93) {$A$}; \node at (2.23,1.93) {$B$};
  \node at (0,-2.46) {$C$};
  \node[anchor=north east] at (3.05,2.1) {$U$};
\end{tikzpicture}
\end{center}

\textbf{Rychlé ověření:} Co znamená být uvnitř obou kruhů?
Kde leží doplněk $A'_U$? Proč každý prvek píšeme do diagramu jen jednou?
Kde v třímnožinovém diagramu leží prvky patřící právě dvěma množinám?

\poznamka{Při dalším setkání použij diagramy k vysvětlení vlastností
množinových operací, například komutativity a De Morganových zákonů.
Prázdné diagramy lze promítnout i vytisknout jako pracovní podklad.}

\end{document}
